\section{Research Question}
\label{sec:checkpoints-question}

\subsection{Problem Statement}
\label{sec:checkpoints-question-problem}

Research programs in process intelligence face a fundamental epistemic problem: the claim that a body
of research is \emph{complete enough to authorize downstream implementation} is itself a claim that
requires evidence. Without a formal verdict mechanism, completeness is asserted rather than
demonstrated. Implementation proceeds on assumption, and the assumption is never falsifiable because
no threshold was declared before the fact.

The checkpoint ledger addresses this problem directly. The RESEARCH\_CRITERIA.md document encodes
the gate criteria \emph{before} the ALIVE verdict is sought. Each criterion names a directory, a
minimum file count, and a verification command. The threshold for ALIVE certification is stated
explicitly: all criteria must be met simultaneously. A single unmet threshold produces PARTIAL. This
means the verdict is falsifiable: it can be re-run, the counts can be checked independently, and any
discrepancy between the stated verdict and the actual artifact inventory is a defect, not a discrepancy.

The problem being solved is not administrative bookkeeping. It is the problem of process-mining
conformance applied reflexively: how can a research program certify that its own knowledge-production
process is sound?

\subsection{Old-Regime Assumption Challenged}
\label{sec:checkpoints-question-old-regime}

The old-regime assumption is that research readiness is a qualitative judgment: a principal investigator
reviews the body of work, concludes that it is \textquotedblleft comprehensive enough,\textquotedblright{}
and authorizes downstream work. No threshold is stated before the fact. No audit trail records which
criteria were evaluated. No distinction is drawn between a complete body of evidence and a
nearly-complete body of evidence with open gaps. Downstream work proceeds in both cases.

This assumption is challenged on two grounds. First, it violates the process-mining principle that
model-versus-log mismatch is a defect, not a discrepancy. If the declared research model requires
fifteen doctrine files and the actual inventory contains twelve, that is a conformance failure, not a
judgment call. Second, it produces false ALIVE verdicts---authorizations that are issued before the
authorizing evidence exists. The RESEARCH\_CRITERIA.md encodes the explicit rule: \textquotedblleft
False ALIVE is breach. PARTIAL is honorable.\textquotedblright{}

\subsection{Hypothesis}
\label{sec:checkpoints-question-hypothesis}

The hypothesis is that a formal, threshold-based, pre-declared gate system---with immutable verdict
documents, cryptographic sealing, and explicit authorization lists---can replace qualitative
research readiness judgment with a falsifiable, auditable, process-conformant certification mechanism.
Specifically: (1) every ALIVE verdict will be traceable to a gate assessment document recording exact
artifact counts against minimum thresholds; (2) every authorization list will name only workflows
whose prerequisites are verified; (3) PARTIAL verdicts will carry explicit bills of materials for the
next transition rather than being treated as failure states; and (4) the Petri net model of the
research lifecycle will reach a sound final marking with no residual tokens in intermediate places.

\subsection{Success Criteria}
\label{sec:checkpoints-question-success}

Success is operationally defined by the gate criteria in RESEARCH\_CRITERIA.md and confirmed by the
ALIVE\_GATE\_ASSESSMENT document. The primary success criterion is 12 of 12 gate criteria met, as
recorded in ALIVE\_GATE\_ASSESSMENT.md (2026-05-31): doctrine 30 of 15 minimum, standards 52 of 10
minimum, papers 15 of 7 minimum, PM4Py 14 of 4 minimum, wasm4pm 20 of 4 minimum, wasm4pm-compat 16
of 4 minimum, lifecycle 42 of 8 minimum, comparisons 11 of 4 minimum, crosswalks 8 of 3 minimum,
M\&A 40 of 6 minimum, adversarial 5 of 3 minimum, gaps 4 of 2 minimum. The secondary success criterion
is that all downstream authorizations issued by ALIVE verdicts name only workflows whose
prerequisites are verified---that no manufacturing work is authorized without a corresponding receipt.
The tertiary success criterion is that the Petri net snapshot records a sound final marking:
[Decommissioning]:1, all intermediate places at zero, confirming deadlock-free process completion.
All three criteria are met by the evidence in the ledger as of 2026-06-01.
